\documentclass[10pt,a4paper]{article}

\usepackage[margin=2.0cm]{geometry}
\usepackage[T1]{fontenc}
\usepackage{lmodern}
\usepackage{microtype}
\usepackage{xcolor}
\usepackage{booktabs}
\usepackage{enumitem}
\usepackage{titlesec}
\usepackage{listings}
\usepackage{amsmath}
\usepackage[hidelinks]{hyperref}

\definecolor{rule}{gray}{0.45}
\definecolor{shade}{gray}{0.94}

\titleformat{\section}{\large\bfseries}{\thesection}{0.6em}{}
\titleformat{\subsection}{\normalsize\bfseries}{\thesubsection}{0.6em}{}
\titlespacing*{\section}{0pt}{1.1ex plus .2ex}{0.6ex}

\lstdefinelanguage{Lean4}{
  morekeywords={def,theorem,structure,class,inductive,instance,namespace,end,
    import,open,where,deriving,axiom,example,by,intro,cases,rw,simp,exact,
    unfold,omega,decide,fun,let,in,if,then,else,do,match,with,return,pure,
    private,partial,Type,Prop,Nat,Bool,UInt8,Array,List,Option,Except},
  sensitive=true,
  morecomment=[l]{--},
  morecomment=[s]{/-}{-/},
  morestring=[b]",
}
\lstset{
  language=Lean4,
  basicstyle=\ttfamily\footnotesize,
  columns=fullflexible,
  keepspaces=true,
  breaklines=true,
  frame=single,
  rulecolor=\color{rule},
  framesep=3pt,
  xleftmargin=4pt,
  xrightmargin=4pt,
  aboveskip=0.6em,
  belowskip=0.3em,
  commentstyle=\itshape\color{black!65},
  keywordstyle=\bfseries,
  showstringspaces=false,
  extendedchars=true,
  literate=
    {×}{{$\times$}}1 {→}{{$\to$}}1 {≤}{{$\le$}}1 {∀}{{$\forall$}}1
    {¬}{{$\neg$}}1 {∧}{{$\wedge$}}1 {∨}{{$\vee$}}1 {⟨}{{$\langle$}}1
    {⟩}{{$\rangle$}}1 {α}{{$\alpha$}}1 {β}{{$\beta$}}1 {ρ}{{$\rho$}}1
    {Γ}{{$\Gamma$}}1 {Σ}{{$\Sigma$}}1 {∈}{{$\in$}}1 {≃}{{$\simeq$}}1
    {↦}{{$\mapsto$}}1 {←}{{$\leftarrow$}}1 {·}{{$\cdot$}}1 {∃}{{$\exists$}}1
    {∉}{{$\notin$}}1 {≥}{{$\ge$}}1
}

\newcommand{\file}[1]{\texttt{#1}}
\newcommand{\lean}[1]{\texttt{#1}}

\title{The Capability Ledger, From the Ground Up\\{\large A tutorial on the
\file{Cap/} subsystem}}
\author{}
\date{}

\begin{document}
\maketitle

\begin{abstract}
\noindent This is a tutorial on the executable half of the repository: how
Rust-like ownership rules are enforced in Lean \emph{from outside} the type
system, by a checked ledger rather than by types. Every declaration named here
exists in the tree and is kernel-checked; \texttt{scripts/check\_all.sh}
rebuilds and re-checks everything, and the axiom audit admits only
\texttt{propext}, \texttt{Classical.choice} and \texttt{Quot.sound}. The code
listings are the real declarations, shortened only by elision of proof bodies.
\end{abstract}

\section{The problem}

Lean's type system has no linearity. A program holding a mutable array can keep
a second reference and observe---or break---an in-place update. Rust solves this
with ownership and borrowing, enforced by the borrow checker. The question here
is whether the same rules can be enforced by a \emph{checked data structure}:

\begin{quote}
A ledger owns all authority. Capabilities are inert names. Every use is checked
against the current state.
\end{quote}

A capability is just a number. It does nothing by itself; all power comes from
the ledger saying, at the moment of use, ``this capability may write these
bytes''. Copying a name is free and grants nothing, because copying is not using.

The stack, bottom to top, each layer trusting only the ones below it:

\begin{center}\footnotesize\ttfamily
\begin{tabular}{ll}
\toprule
algorithms as clients & \file{Cap/UFRun.lean} \\
sealed certificates & \file{Cap/Runner.lean} \\
unified checker/interpreter & \file{Cap/Unified.lean} \\
suspension (loans) & \file{Cap/Suspend.lean} \\
the world & \file{Cap/World.lean}, \file{Cap/WorldInvariant.lean} \\
context and invariants & \file{Cap/Context.lean}, \file{Cap/Invariant.lean} \\
joins & \file{Cap/Join.lean} \\
loan bookkeeping & \file{Cap/Loans.lean} \\
the ledger & \file{Cap/Ledger.lean} \\
shares & \file{Cap/Share.lean} \\
footprints & \file{Cap/Footprint.lean} \\
identities & \file{Cap/Id.lean} \\
\bottomrule
\end{tabular}
\end{center}

\section{Names are not pointers}

Everything the ledger tracks has its own identity type (\file{Cap/Id.lean}):

\begin{lstlisting}
structure AllocId where
  val : Nat
deriving DecidableEq, Repr

structure CapId where
  val : Nat
deriving DecidableEq, Repr
-- and likewise LoanId, LifetimeId
\end{lstlisting}

They are distinct tagged naturals --- you write \lean{⟨3⟩} --- and being distinct
\emph{types} means the checker can never confuse an allocation for a capability.

\section{Shares: who may do what}

A capability's \emph{share} of a resource is a fraction $q/D$: a numerator $q$
over a context-wide denominator $D$. The whole resource is $D/D$; a half is
$D/2$ over $D$. \file{Cap/Share.lean} enforces three rules:

\begin{itemize}[leftmargin=1.4em]
\item a split produces shares that add up to exactly what was there
  (\lean{Splits});
\item only the \emph{whole} share may write (\lean{Writable});
\item so two partial shares can never both write (\lean{no\_two\_writes}).
\end{itemize}

That is the whole fractional-uniqueness story: writing requires $q = D$, and the
ledger's accounting makes $q = D$ achievable by at most one capability at a
time.

\section{Footprints: which bytes}

A footprint is a half-open interval $[lo, hi)$ of an allocation's bytes
(\file{Cap/Footprint.lean}). \lean{splitAt} cuts an interval in two; the proofs
say the pieces are disjoint, adjacent and cover the original. A capability's
authority is a share \emph{and} a footprint together: this much of these bytes.

\section{The ledger: trees of authority}

An \lean{Owner} (one per allocation) holds a tree of leaves
(\file{Cap/Ledger.lean}). A \lean{Leaf} is a capability with its footprint, its
share numerator, a ceiling and a state:

\begin{lstlisting}
inductive CapState where
  | active
  | retired
  | consumed

structure Leaf where
  id : CapId
  alloc : AllocId
  lo : Nat
  hi : Nat
  ordered : lo ≤ hi
  q : Nat
  pos : 0 < q
  ceiling : Ceiling
  state : CapState

inductive Tree where
  | leaf (c : Leaf)
  | frac (id : CapId) (l r : Tree)      -- children split the share, same bytes
  | spatial (id : CapId) (l r : Tree)   -- children split the byte range
\end{lstlisting}

The measure that makes it all cohere is the \emph{budget}: at byte $x$, the sum
of shares of all leaves covering $x$.

\begin{lstlisting}
def Leaf.budget (c : Leaf) (x : Nat) : Nat :=
  if c.state = .consumed then 0
  else if c.lo ≤ x ∧ x < c.hi then c.q else 0
\end{lstlisting}

The design's central invariant is coverage (\file{Cap/Ledger.lean}):

\begin{lstlisting}
def Owner.Covered (o : Owner) (D : Nat) : Prop :=
  ∀ x, x < o.len → o.root.budget x = D
\end{lstlisting}

At every byte the budget is exactly $D$: nothing unaccounted for, nothing
double-counted. Splits preserve the budget by construction
(\lean{frac\_split\_budget}, \lean{spatial\_split\_budget}); the sibling-only
fractional join merges two leaves back into one of summed share and preserves it
(\lean{joinFraction\_covered}), and is refused if the children carry the same
capability id (\lean{joinChildren?\_alias\_none}).

The write guard is decidable, and provably the same thing as the logical guard:

\begin{lstlisting}
def Owner.writableB (o : Owner) (D : Nat) (id : CapId) (i : Nat) : Bool :=
  match o.root.find? id with
  | none => false
  | some c =>
    decide (c.state = .active ∧ c.q = D ∧ c.ceiling = .readWrite
              ∧ c.lo ≤ i ∧ i < c.hi)

theorem Owner.writableB_sound
    (h : o.writableB D id i = true) : o.Writable D id i
\end{lstlisting}

\lean{writableB\_sound} is small but load-bearing: it is why the executable
check and the logical predicate can never drift apart.

\section{Many owners: the context}

A \lean{Ctx} is a list of owners, the denominator and a fresh-allocation
counter (\file{Cap/Context.lean}):

\begin{lstlisting}
structure Ctx where
  owners : List Owner
  D : Nat
  nextAlloc : Nat

def Ctx.Realizes (C : Ctx) : Prop :=
  (∀ o ∈ C.owners, o.Covered C.D) ∧ (C.owners.map (·.alloc)).Nodup
\end{lstlisting}

\file{Cap/Invariant.lean} strengthens this to \lean{Ctx.WellFormed}:
\lean{Realizes}, plus every owner's leaves inside its own allocation
(\lean{Owner.LeavesWithin}), plus all capability ids across all owners distinct.
The important theorems are the \emph{preservation} ones --- allocation,
disposal, and both splits, dispatched to their owner, all preserve it:

\begin{lstlisting}
theorem Ctx.splitAtOwner_wellFormed
    (hwf : C.WellFormed) (hmem : o ∈ C.owners)
    (hfids : o.FreshIds id1 id2)
    (hfresh1 : id1 ∉ C.ids) (hfresh2 : id2 ∉ C.ids)
    (hs : C.splitAtOwner? o id id1 id2 k = some C') :
    C'.WellFormed
\end{lstlisting}

Freshness is not a detail: it is what stops a new capability from colliding with
an old one and inheriting its authority. The join side,
\lean{Ctx.joinFractionOwner\_wellFormed}, is harder: joining removes two
capability ids and adds one, so its proof needs an exact permutation of the id
list (\lean{Owner.joinFraction?\_ids\_perm}).

\section{The world: owners plus bytes}

A \lean{World} (\file{Cap/World.lean}) is the context, the backends (one byte
buffer per allocation), and the loan bookkeeping:

\begin{lstlisting}
structure World where
  ctx : Ctx
  heaps : AllocId → Option Backend
  loans : LoanCtx
\end{lstlisting}

\file{Cap/WorldInvariant.lean} bundles the world invariant:

\begin{lstlisting}
def World.WellFormed (w : World) : Prop :=
  w.ctx.WellFormed           -- context invariant
  ∧ w.ctx.Functional         -- one owner per allocation
  ∧ w.StorageMatches         -- the found owner's backend has its length
\end{lstlisting}

The workhorse lemma is \lean{World.update\_storage}: replacing one owner by a
same-length result keeps the heap agreement. Every command's preservation proof
ends in it.

\section{The unified language: the checker is the interpreter}

\file{Cap/Unified.lean} gives the command language:

\begin{lstlisting}
inductive UCmd where
  | write (a : AllocId) (id : CapId) (i : Nat) (v : UInt8)
  | splitAt (a : AllocId) (id id1 id2 : CapId) (k : Nat)
  | splitFraction (a : AllocId) (id id1 id2 : CapId) (q1 q2 : Nat)
  | join (a : AllocId) (node newId : CapId)
  | free (a : AllocId) (d : CapId)
  | alloc (n : Nat) (rootId disposal : CapId)
  | beginLoan (a : AllocId) (lender : CapId) (loanId : LoanId)
  | endLoan (a : AllocId) (loanId : LoanId)
  | nop
\end{lstlisting}

\lean{checkWorld} runs a command list and returns either an error or the world
it executed:

\begin{lstlisting}
def checkWorld : World → List UCmd → Except String World
  | w, [] => .ok w
  | w, .nop :: cs => checkWorld w cs
  | w, .write a id i v :: cs =>
    match w.writeChecked a id i v with
    | .ok w' => checkWorld w' cs
    | .error e => .error e
  -- and likewise one branch per command
\end{lstlisting}

One theorem covers all of it:

\begin{lstlisting}
theorem checkWorld_wf :
    checkWorld w cs = .ok w' → w.WellFormed → w'.WellFormed
\end{lstlisting}

Every accepted execution --- allocation, writes, both splits, joins, loans,
disposal --- preserves the world invariant. \lean{checkWorld\_rel} derives the
older \lean{World.Rel} from it, and \lean{checkWorld\_append}
(\file{Cap/Runner.lean}) is what makes sequential composition of certificates
possible.

\subsection{A walk through the lifecycle}

The concrete fixture runs the whole discipline end to end:

\begin{lstlisting}
def uLifecycle : List UCmd :=
  [ .alloc 4 ⟨3⟩ ⟨4⟩             -- fresh 4-byte owner ⟨1⟩; root ⟨3⟩, disposal ⟨4⟩
  , .write ⟨1⟩ ⟨3⟩ 0 7           -- whole-share write at byte 0
  , .splitAt ⟨1⟩ ⟨3⟩ ⟨5⟩ ⟨6⟩ 2   -- ⟨5⟩ covers [0,2), ⟨6⟩ covers [2,4)
  , .write ⟨1⟩ ⟨5⟩ 0 8           -- write through the left half
  , .splitFraction ⟨1⟩ ⟨5⟩ ⟨7⟩ ⟨8⟩ 32 32  -- two half-shares of [0,2)
  , .join ⟨1⟩ ⟨7⟩ ⟨9⟩            -- merge back to a whole share ⟨9⟩
  , .write ⟨1⟩ ⟨9⟩ 1 9           -- writable again: the join restored it
  , .free ⟨1⟩ ⟨4⟩ ]              -- disposal through the recorded capability
\end{lstlisting}

Between the fractional split and the join, a write through \lean{⟨7⟩} would be
\emph{refused} --- a half share cannot write --- and the budget arithmetic is
what makes that decision sound. The evidence is kernel-computed:

\begin{lstlisting}
theorem uLifecycle_ok : (checkWorld uWorld uLifecycle).isOk = true := by decide

theorem uLifecycle_wf :
    ∃ w, checkWorld uWorld uLifecycle = .ok w ∧ w.WellFormed

theorem uLifecycle_snapshot :   -- the untouched 64-byte backend reads back zeros
    uLifecycleSnapshot = some (some (List.replicate 64 (0 : UInt8))) := by decide
\end{lstlisting}

\section{Loans: suspension, not tracking}

Loans reuse the leaf state the write guard already consults
(\file{Cap/Suspend.lean} plus the loan commands in \file{Cap/Unified.lean}).
\lean{beginLoan} retires the lender's leaf; since \lean{writableB} requires
\lean{.active}, a write through the lender is then refused \emph{by the same
guard} --- no parallel machinery:

\begin{lstlisting}
def Owner.suspendLeaf? (o : Owner) (id : CapId) : Option Owner
  -- succeeds iff the leaf is .active; sets it to .retired

theorem Owner.suspendLeaf?_not_writable :
    o.suspendLeaf? id = some o' → ∀ i, o'.writableB D id i = false

theorem Owner.resumeLeaf?_of_suspendLeaf? :
    -- suspend then resume returns the original owner
\end{lstlisting}

The refusal is a theorem with a positive control: the same write is accepted
before the loan (\lean{uLoan\_write\_before}), refused while it is outstanding
(\lean{uLoan\_suspended\_write\_refused}), and accepted again after
\lean{endLoan} (\lean{uLoanRun\_ok}).

Stated limits, on purpose: the borrower is the lender's own id (exclusive
suspension, not a reborrow capability), and the \file{Cap/Policy.lean} liveness
conjunction is a predicate form the checker does not consult.

\section{Certificates and the sealed boundary}

A \lean{Prog pre} (\file{Cap/Runner.lean}) is a command list packaged with the
proof that the checker accepts it --- you cannot construct one without the
derivation:

\begin{lstlisting}
structure Prog (pre : World) where
  post : World
  cmds : List UCmd
  ok : checkWorld pre cmds = .ok post

theorem checkWorld_append :
    checkWorld w cs = .ok w' → checkWorld w (cs ++ ds) = checkWorld w' ds

def Prog.seq (p : Prog pre) (q : Prog p.post) : Prog pre
  -- the second program provably starts in the first's post-world
\end{lstlisting}

\lean{ClosedProgram} is the sealed form: commands accepted from the canonical
empty world, with the snapshot as the only observation:

\begin{lstlisting}
def World.emptyWorld : World
  -- no owners, no backends, one live root lifetime

structure ClosedProgram where
  cmds : List UCmd
  accepted : (checkWorld World.emptyWorld cmds).isOk = true

theorem closed64_snapshot :   -- allocate 64, fill both halves, read back
    (checkWorld World.emptyWorld closed64.cmds).toOption.map
      (fun w => (w.heaps ⟨0⟩).map (·.buf))
      = some (some (List.replicate 32 1 ++ List.replicate 32 2)) := by decide
\end{lstlisting}

The residual honesty: Lean has no module opacity, so \lean{World} itself remains
public --- the seal is by discipline, not by the module system.

\section{Algorithms as clients}

Union-find runs over the checked world (\file{Cap/UFRun.lean}). Reads are
observations of the pure backends --- the discipline constrains writes --- and
\lean{union} issues its writes through \lean{World.writeChecked}:

\begin{lstlisting}
def find (w : World) : Nat → Nat → Option Nat
  | 0, _ => none
  | fuel + 1, i =>
    match readParent w i with
    | none => none
    | some p => if p = i then some i else find w fuel p

def union (w : World) (i j : Nat) : Except String World :=
  match find w 8 i, find w 8 j with
  | some ri, some rj =>
    if ri = rj then .ok w
    else w.writeChecked ⟨0⟩ ⟨1⟩ ri rj.toUInt8
  | _, _ => .error "union: find failed"

theorem union_wf : union w i j = .ok w' → w.WellFormed → w'.WellFormed
\end{lstlisting}

The concrete trace (attach 1 under 0, 3 under 2, 2 under 0) is accepted, its
final world well formed, and its backend parent array is exactly the functional
model's \lean{[0,0,0,2]}, with every pair of the four nodes \lean{Same} in the
partition relation:

\begin{lstlisting}
theorem ufTrace_partition :
    (∃ w', ufTrace ufWorld = .ok w' ∧ w'.WellFormed ∧
      (w'.heaps ⟨0⟩).map (fun b => (b.buf.take 4).map UInt8.toNat)
        = some [0, 0, 0, 2]) ∧
    ∀ x y, x < 4 → y < 4 → UF.Same [0, 0, 0, 2] x y
\end{lstlisting}

The reachability and one-step-compression facts the client would preserve live
in \file{Cap/UFPartition.lean} and \file{Cap/UFCompress.lean}
(\lean{compress\_same\_sound}: compression from a reachable root creates no new
class pairs).

\section{Two backends, one model}

The pure \lean{Backend} (a byte list with a length proof) is the specification;
\lean{BackendArr} (\file{Cap/BackendArr.lean}) is the same interface over Lean's
\lean{ByteArray}:

\begin{lstlisting}
structure BackendArr where
  arr : ByteArray
  len : Nat
  hlen : arr.size = len

def BackendArr.toBackend (x : BackendArr) : Backend :=
  { buf := x.arr.data.toList, len := x.len, ... }

theorem BackendArr.write_refines :
    (x.write i v).map BackendArr.toBackend = Backend.write x.toBackend i v

theorem ba64_refines : ba64.map BackendArr.toBackend = pure64
theorem ba64_snapshot :   -- the byte-array run lands on the reference bytes
    (ba64.map BackendArr.toBackend).map Backend.buf
      = some (List.replicate 32 1 ++ List.replicate 32 2)
\end{lstlisting}

The tempting next sentence --- ``the byte array is updated in place'' --- is
exactly the sentence this repository refuses to write here: Lean's runtime
\emph{may} mutate a uniquely referenced \lean{ByteArray}, but that is not
observable from inside, so nothing in \file{Cap/} claims it. Observing and
certifying in-place update is the \file{Ownership/} side's subject. The timing
observation (\file{bench/CapBATiming.lean}, one interpreted run, agreeing
checksums) is recorded with the same caveat:

\begin{center}
\begin{tabular}{lr}
\toprule
route & ns / repetition \\
\midrule
pure list backend & 73{,}959 \\
byte-array backend & 47{,}688 \\
reference heap & 51{,}956 \\
\bottomrule
\end{tabular}
\end{center}

\section{What is not claimed}

\begin{itemize}[leftmargin=1.4em]
\item No concurrency, no FFI, no real runtime.
\item No read discipline: reads are free observations.
\item Loans are exclusive suspension; no reborrows, no nested lifetimes in the
  executable language.
\item The sealed boundary is not module-enforced.
\item No in-place-update or compiled-performance claim anywhere in \file{Cap/}.
\item Union-find correctness is the concrete four-node instance, not the
  general theorem.
\end{itemize}

\section{The research-scale remainder}

\begin{enumerate}[leftmargin=1.6em]
\item \textbf{A read discipline} --- fractional read permissions through the
  guard, and a second, non-suspending loan mode.
\item \textbf{General union-find correctness} --- termination for compression
  (a rank potential) plus ownership accounting of the write trace: the first
  theorem that would use the discipline to prove an algorithm right, not just
  safe.
\item \textbf{The in-place observation problem} --- connecting ledger
  exclusivity to Lean's runtime reference counts; the \file{Ownership/} side
  carries this.
\item \textbf{A genuine sealed boundary} --- an opaque runtime or a
  verified-backend embedding; a policy decision about axioms, not just a proof.
\item \textbf{Compiled measurements} --- a native binary with a defensible
  methodology; only interesting once the previous item moves.
\item \textbf{Full lifecycle bookkeeping} --- non-root lifetimes, distinct
  borrowers, disposal refused while any loan is outstanding at any owner.
\end{enumerate}

\section{Reproduce everything}

\begin{lstlisting}[language={},frame=none,backgroundcolor=\color{shade}]
nix develop -c lake build Cap                -- the whole subsystem
nix develop -c bash scripts/check_all.sh     -- builds + axioms + census + docs
lake env lean bench/CapBATiming.lean         -- the three-backend observation
\end{lstlisting}

\end{document}
